\begin{table}[t]
		\caption{与在线建图方法的性能比较。结果在 $[1.0m, 1.5m, 2.0m]$ 阈值下评估。$\ddag$：通过官方开源代码复现。$\dag$：基于\model 中规划模块的需求，将\textit{边界}进一步划分为\textit{道路段}和\textit{车道}分别评估。$*$：骨干网络和稀疏感知模块的成本。-R：带雷达点云输入。}
	        \vspace{-10pt}
	        % \scriptsize
	        \tiny
		\begin{center}
			% \resizebox{1.0\textwidth}{!}
	            {
				\begin{tabular}{l|p{1.3cm}<{\centering}p{1.3cm}<{\centering}p{2cm}<{\centering}p{1cm}<{\centering}p{3.0cm}<{\centering}}
					\toprule
					方法 & \textbf{分隔线}$\uparrow$& \textbf{人行横道}$\uparrow$ & \textbf{边界}$\uparrow$ & mAP$\uparrow$ & 训练内存(GB)$\downarrow$\\
	                \midrule
	                \multicolumn{6}{l}{\textit{纯在线建图方法}}\\
	                \midrule
					HDMapNet$^\ddag$ \cite{li2022hdmapnet} & 18.0 & 23.1 & 33.5 & 24.9  & 21.1\\
					VectorMapNet$^\ddag$ \cite{liu2023vectormapnet}  & 39.8 &  38.7 & 19.9 & 32.8 &16.4\\
					MapTR$^\ddag$ \cite{liao2022maptr}  & 17.9 &  36.9 & 1.8 & 18.9 & 20.0\\
				  StreamMapNet$^\ddag$ \cite{yuan2024streammapnet} & 55.5 &  59.7 & 46.7 & 54.0 &23.2\\
					\midrule
	                \multicolumn{6}{l}{\textit{端到端多目标跟踪方法}}\\
	                \midrule
					\rowcolor[gray]{.9} {SparseAD-B}  & {37.3}& {45.2} & {29.2/25}$^\dag$ & {34.2} & \textbf{6.9}$^*$\\
	                    \rowcolor[gray]{.9} {SparseAD-BR}  & {38.4}& {44.5} & {29.4/24.7}$^\dag$ & {34.2} & 7.1$^*$\\
	                    \rowcolor[gray]{.9} {SparseAD-L}  & {37.6}& {46.7} & {30/24.8}$^\dag$ & {34.8} & 19.3$^*$\\
	                \bottomrule
				\end{tabular}
			}
		\end{center}
		\label{tab:map-nuscenes-val}
	         \vspace{-10pt}
	\end{table}